\SourceLeaf{35}

Es wird dann folglich der Werth des Bruchs $\dfrac{P_\theta}{P_{\theta+1}}$ zu klein
gefunden, wenn man im Zähler die beiden ersten Glieder, und im Nenner nur das erste in
Rechnung bringt, oder mit anderen Worten, der wahre Werth jenes Bruches ist grösser als
$(2\theta+1) . \dfrac{n}{mm}\left\{1-\dfrac{1}{2(2\theta+1)}\dfrac{mm}{nn}\right\}$ oder
$\dfrac{n}{mm}(2\theta+1)-\dfrac{1}{2n}$. 
Es folgt hieraus, dass $P_\theta, P_{\theta+1}, P_{\theta+2}$ u.\,s.\,w. eine stets
abnehmende unendliche Reihe von Grössen bilden, wenn nur $2\theta+1$ nicht kleiner ist als
$\dfrac{mm}{n}+\dfrac{mm}{2nn}$. Es ist also unmöglich, dass sämmtliche Glieder der Reihe
$P, P_1, P_2, P_3$ u.\,s.\,w. ganzen Zahlen proportional seien. Diess würde aber
nothwendig der Fall sein, wenn man annimmt, dass $m$ und $n$ ganze Zahlen sind und
zugleich die beiden ersten Glieder jener Reihe $P, P_1$ zu einander in einem rationalen
Verhältniss stehen. Man hat nemlich:

\[
\begin{aligned}
P_2 &= nP_1 - mP\\
P_3 &= 3nP_2 - mmP_1\\
P_4 &= 5nP_3 - mmP_2\\
P_5 &= 7nP_4 - mmP_3\\
&\text{u.\,s.\,w.,}
\end{aligned}
\]
\noindent allgemein, von der zweiten Formel an,

\[
P_{\theta+2} = (2\theta+1)nP_{\theta+1} - mmP_\theta.
\]

Offenbar kann man in diesen Formeln anstatt der Grössen $P, P_1, P_2$ u.\,s.\,w. auch
ihnen proportionale setzen. Wären aber von diesen die zwei ersten ganze Zahlen, so würden
auch alle folgenden ganze Zahlen sein, was, wie eben gezeigt wurde, unmöglich ist.

Es kann folglich, wenn $m, n$ ganze Zahlen bedeuten, der Werth des Bruchs
$\dfrac{P_1}{P}$, welcher nichts anderes ist als die Tangente des Bogens
$\dfrac{m}{n}$, nicht rational sein.
